\section{Selected Inference Algorithms for MixDom}

\subsection{Fast Statistical Alignment (FSA)}
\label{sec:fsa}

Given a set of sequences and a phylogenetic tree with branch-specific pair HMMs
(TKF92, MixDom, or distilled order-1 transducers),
we construct a multiple sequence alignment using the
\emph{sequence annealing} approach of \cite{BradleyEtAl2009},
to which we refer the reader for a full description of the algorithm.

Briefly, the method proceeds as follows.
For each pair of sequences $(x,y)$ in a selected subset
(either all $\binom{N}{2}$ pairs or an $O(N \log N)$ Erd\H{o}s--R\'enyi sample),
we compute pairwise residue alignment posteriors $P(x_i \sim y_j)$
by running the Forward-Backward algorithm on the pair HMM at an
optimized evolutionary time $\hat{\tau}$.
The time $\hat{\tau}$ is found by Newton--Raphson optimization
of the expected log-likelihood (the ``NR step''):
\begin{align}
\hat{\tau} &= \operatorname*{argmax}_\tau \; \mathbb{E}_{P(\pi | x,y,\tau_0)} \bigl[ \log P(x,y,\pi \mid \tau) \bigr]
\label{eq:fsa-nr}
\end{align}
where $\pi$ ranges over alignment paths and $\tau_0$ is an initial estimate.
This expectation is computed from Forward-Backward expected counts
at $\tau_0$, and typically converges in 3--5 Newton steps.
(This time-maximization differs slightly from the approach of \cite{BradleyEtAl2009}
which attempts to optimize all model parameters via unregularized EM for every pair,
and consequently must terminate the EM recursion early to avoid instability.)

The pairwise posteriors are then assembled into a multiple alignment
by the greedy sequence annealing procedure of \cite{BradleyEtAl2009},
which iteratively merges alignment columns to maximize a sum-of-pairs
posterior objective.


\subsection{Beam Search Ancestral Sequence Reconstruction (BeamASR)}
\label{sec:beam-ancestor}

We now describe an alternative progressive reconstruction method
that finds the maximum-likelihood ancestral sequence at each internal node
by beam search, without materializing the full composite automaton.

At each internal node $v$ with children $l,r$ and observed descendant sequences
$c_l, c_r$, we seek
\begin{align}
\hat{a}_v &= \operatorname*{argmax}_{a} \bigl[ \log P(a, c_l \mid B_l) + \log P(a, c_r \mid B_r) - \log P(a \mid R) \bigr]
\label{eq:beam-ancestor}
\end{align}
where $P(a, c_k \mid B_k)$ is the pair HMM forward probability on branch $k$
and $P(a \mid R)$ is the singlet probability under the root generator,
subtracted to avoid double-counting the prior on~$a$.

\subsubsection{Incremental forward profiles}

Since the branches are conditionally independent given the ancestor,
we can evaluate~(\ref{eq:beam-ancestor}) by maintaining
\emph{incremental forward profiles}: for each branch $k$, a 1D forward table
$F_k[i, q]$ giving the log-probability that descendant positions $1,\ldots,i$
have been emitted and the branch machine is in state $q$,
given ancestor positions $1,\ldots,j$ processed so far.

Each ancestor character extends both profiles independently in $O(L_k)$ time
per branch, where $L_k = |c_k|$.

\subsubsection{Beam search}

The ancestor sequence $\hat{a}_v$ is built left-to-right by beam search.
At each position $j$, the beam maintains $B$ candidate partial ancestors.
For each candidate and each alphabet character $\sigma$:
\begin{enumerate}
\item Extend both branch profiles by one ancestor character $\sigma$,
  comprising a \emph{match/delete phase} (the ancestor emits $\sigma$,
  descendant positions advance via M or D transitions)
  and an \emph{insertion phase} (descendant-only insertions following the
  ancestor emission).
\item Update the singlet forward score for $\sigma$.
\item Score the extension:
  $\Delta(j, \sigma) = \Delta F_{l} + \Delta F_{r} - \Delta_{\mathrm{singlet}}$.
\end{enumerate}
The top $B$ extensions (by cumulative score) are retained.
Total cost per node is $O(K \cdot B \cdot A \cdot (L_l + L_r))$
where $K = |\hat{a}_v|$ and $A$ is the alphabet size.

\subsubsection{Insertion phase via associative scan}

The insertion recurrence within each branch profile has the form
\begin{align}
x_{i+1} &= \operatorname{logsumexp}\bigl(A_{II} \, x_i,\; b_i\bigr) + e_i
\label{eq:insert-recurrence}
\end{align}
where $A_{II}$ is the I-to-I log-transition submatrix, $b_i$ collects
transitions into insertion states from M and D, and $e_i$ is the emission score.
This is a log-semiring affine recurrence, parallelizable via
an associative scan with operator
\[
(A_1, b_1) \oplus (A_2, b_2) = (A_2 \otimes A_1,\; \operatorname{logsumexp}(A_2 \, b_1,\; b_2))
\]
where $\otimes$ denotes log-semiring matrix multiplication.
This reduces the insertion phase from $O(L)$ sequential depth to $O(\log L)$.

\subsubsection{Supported model types}

The beam search interface is generic over the pair HMM used on each branch:
\begin{enumerate}
\item \textbf{TKF92} --- order-0, 5-state pair HMM (the standard model).
\item \textbf{MixDom} --- the full latent-state pair HMM (Section~\ref{sec:nestpairhmm})
  with $2 + 5NK$ states; latent-state correlations are marginalized
  in the forward pass without distillation.
%\item \textbf{Distilled order-1 transducer} --- an order-1 context-dependent
%  transducer (Section~\ref{sec:order1trans}), expanded to a standard pair HMM
%  with $O(A)$ states per structural state.
\end{enumerate}


%\subsection{Variational ancestral reconstruction (VarAnc)}
% \label{sec:varanc}
% 
% Given an MSA with known gap structure---i.e., for each internal node $v$ and MSA column $c$,
% we know whether $v$ is present (ungapped) or absent (gapped)---but unknown ancestral residue
% identities, we optimize a variational lower bound (ELBO) on the marginal likelihood.
% Our approach is a ``product of trees'' structured variational
% approximation~\cite{JojicEtAl2004}, adapted to the order-1 WFST framework.
% 
% \subsubsection{Model structure}
% 
% The joint distribution over observed leaf sequences $\mathbf{y}$
% and ancestral sequences $\mathbf{a}$ factorizes over tree edges.
% With the alignment structure fixed, the log-likelihood decomposes into
% contributions from consecutive alignment-column pairs along each edge $e = (u,v)$:
% \begin{align}
% \log P(\mathbf{y}, \mathbf{a})
%  &= \sum_{e \in \mathrm{edges}} \sum_{k=1}^{K_e}
%     \log w_e\!\bigl(\mathrm{type}_{k-1},\, \mathrm{type}_k,\, \mathrm{ctx}_{k-1},\, \mathrm{ctx}_k\bigr)
% \label{eq:lik-edges}
% \end{align}
% where $w_e$ is the order-1 WFST transition weight on edge $e$
% (depending on the previous and current column types and context characters).
% The first and last columns contribute start ($\sta$) and end ($\fin$) transition factors.
% 
% The key computational challenge is that the order-1 transitions couple
% adjacent alignment columns: the WFST weight at column $c$ depends on the
% ancestral and descendant characters at both column $c$ and its predecessor.
% This makes the full graphical model a tree (phylogeny) $\times$ chain (sequence)
% grid with undirected cycles, and exact inference is intractable.
% 
% \subsubsection{Product-of-trees variational distribution}
% 
% Following~\cite{JojicEtAl2004}, we factorize the variational posterior
% over MSA columns while retaining the full tree joint at each column:
% \begin{align}
% q(\mathbf{a}) &= \prod_{c=1}^{L} q_c\!\bigl(\mathbf{h}_c\bigr)
% \label{eq:var-q-pot}
% \end{align}
% where $\mathbf{h}_c = \{a_{v,c} : v \in \mathrm{ungapped}(c)\}$ denotes the set of
% all ancestral characters at MSA column $c$,
% and each $q_c$ is a tree-structured distribution over the characters at the
% internal nodes that are ungapped at column $c$.
% Leaf characters are observed (clamped).
% 
% This approximation decouples columns in $q$ but preserves the
% full parent-child correlation structure within each column.
% In particular, the pairwise joint marginal
% $q_c^{(u,v)}(a_u, a_v)$ on each tree edge is available exactly from
% the Felsenstein peeling/unpeeling pass used to represent $q_c$
% (Section~\ref{sec:felsenstein-update}).
% 
% \subsubsection{Free energy and ELBO}
% 
% The evidence lower bound is
% \begin{align}
% \mathcal{L}(q) &= \mathbb{E}_q\bigl[\log P(\mathbf{y}, \mathbf{a})\bigr] + H(q)
%   \;\geq\; \log P(\mathbf{y})
% \label{eq:elbo}
% \end{align}
% where $H(q) = -\sum_c \sum_{\mathbf{h}_c} q_c(\mathbf{h}_c) \log q_c(\mathbf{h}_c)$
% is the entropy of the product-of-trees distribution.
% Since each $q_c$ is tree-structured, its entropy decomposes as
% \begin{align}
% H(q_c) &= \sum_{v \in \mathrm{ungapped}(c)} H\!\bigl(q_c^{(v)}\bigr)
%   - \sum_{e = (u,v)} \mathrm{MI}_c(u, v)
% \label{eq:tree-entropy}
% \end{align}
% where $H(q_c^{(v)})$ is the marginal entropy at node $v$ and
% $\mathrm{MI}_c(u,v)$ is the mutual information between parent $u$ and child $v$
% at column $c$, both computable from the Felsenstein marginals.

\subsubsection{Potentials from neighboring columns}

The inter-column coupling enters through the order-1 WFST transitions.
For each edge $e = (u,v)$ and MSA column $c$ where $e$ has an event of type $\tau_c$,
let $c^-$ denote the predecessor column (the previous column where $e$ had an event)
and $c^+$ the successor column.
The \emph{potential} at column $c$ for edge $e$ receives two contributions:

\paragraph{As-child term (from $c^-$).}
The transition from column $c^-$ to $c$ on edge $e$ depends on the
characters at both columns. Using the pairwise marginal from $q_{c^-}$:
\begin{align}
\log \phi_e^{\mathrm{child}}(a_{u,c}, a_{v,c})
 &= \sum_{a', b'} q_{c^-}^{(u,v)}(a', b')\;
    \log w_e(\tau_{c^-\!}, \tau_c, a', b', a_{u,c}, a_{v,c})
\label{eq:pot-child}
\end{align}
where the sum over $(a', b')$ uses the \emph{joint} pairwise marginal
$q_{c^-}^{(u,v)}(a', b')$---not the product of independent marginals.
This is the key advantage over mean-field:
the within-tree parent-child correlation at the predecessor column is
preserved exactly.

\paragraph{As-parent term (from $c^+$).}
Symmetrically, column $c$ acts as the predecessor for column $c^+$:
\begin{align}
\log \phi_e^{\mathrm{parent}}(a_{u,c}, a_{v,c})
 &= \sum_{a'', b''} q_{c^+}^{(u,v)}(a'', b'')\;
    \log w_e(\tau_c, \tau_{c^+\!}, a_{u,c}, a_{v,c}, a'', b'')
\label{eq:pot-parent}
\end{align}

For insert transitions (only the descendant is present at $c$),
the potential reduces to a per-node function $\psi_v(a_{v,c})$.
For delete transitions (only the ancestor is present), it becomes $\psi_u(a_{u,c})$.
For match transitions, it contributes a per-edge potential $\phi_e(a_{u,c}, a_{v,c})$.
The start and end transitions contribute analogous per-node or per-edge terms.

\subsubsection{Felsenstein coordinate ascent}
\label{sec:felsenstein-update}

Each coordinate ascent step updates $q_c$ for a single MSA column $c$,
holding all other columns fixed.
We accumulate, for each edge $e$ and node $v$ in the tree at column $c$:
\begin{itemize}
\item Per-edge log-potentials
  $\log \phi_e(a_u, a_v) = \log \phi_e^{\mathrm{child}} + \log \phi_e^{\mathrm{parent}}$
  (for match transitions where both endpoints are present).
\item Per-node log-potentials
  $\log \psi_v(a_v)$ (from insert/delete transitions on incident edges,
  plus the root prior $\log \pi(a)$ at the root node).
\end{itemize}

The optimal $q_c$, given the potentials, is the Gibbs distribution on the
tree at column $c$:
\begin{align}
q_c(\mathbf{h}_c) &\propto \prod_v \psi_v(a_v) \prod_{e=(u,v)} \phi_e(a_u, a_v)
  \prod_{\ell \in \mathrm{leaves}} \delta(a_\ell = y_\ell)
\label{eq:q-gibbs}
\end{align}

Since this is a tree-structured MRF, the normalizing constant and all
node and edge marginals can be computed \emph{exactly} by Felsenstein
peeling (postorder) and unpeeling (preorder) in $O(|E| \cdot |\alphabet|^2)$ time.

\paragraph{Peeling (postorder).}
For each node $v$ in postorder, compute the conditional likelihood:
\begin{align}
\mathrm{CL}_v(a) &= \psi_v(a) \prod_{\mathrm{children}\; c}
  \Bigl[\sum_{a_c} \phi_{(v,c)}(a, a_c)\; \mathrm{CL}_c(a_c)\Bigr]
\label{eq:peel}
\end{align}
with $\mathrm{CL}_\ell(a) = \delta(a = y_\ell)$ for observed leaves.
The log-partition function is $\log Z_c = \log \sum_a \pi(a)\, \mathrm{CL}_{\mathrm{root}}(a)$.

\paragraph{Unpeeling (preorder).}
Propagate top-down to obtain the posterior marginal at each node:
\begin{align}
q_c^{(v)}(a) &\propto \mathrm{CL}_v(a) \cdot \mathrm{msg}_{\mathrm{parent} \to v}(a)
\label{eq:unpeel-marginal}
\end{align}
and the pairwise marginal on each edge:
\begin{align}
q_c^{(u,v)}(a_u, a_v)
 &\propto \mathrm{msg}_{\mathrm{above}\,u}(a_u) \cdot \phi_{(u,v)}(a_u, a_v) \cdot \mathrm{CL}_v(a_v)
\label{eq:unpeel-pair}
\end{align}
where the ``message from above $u$'' combines the top-down message to $u$
with $u$'s conditional likelihood excluding child $v$.

\paragraph{Sweep.}
One iteration sweeps through all MSA columns $c = 1, \ldots, L$:
for each column, recompute the potentials from the current neighbor marginals,
run peeling/unpeeling, and store the updated node and edge marginals.
The sweep order is left-to-right; the ``as-child'' potentials use the
just-updated predecessor marginals, while the ``as-parent'' potentials
use stale successor marginals from the previous iteration.

\subsubsection{Properties}

The product-of-trees approximation enjoys the same monotonic convergence
guarantee as mean-field coordinate ascent (each column update minimizes
the free energy in its coordinate), with the additional guarantee that
the ELBO is at least as tight as the fully-factored mean-field bound.
This follows because the product-of-trees family \emph{contains} the
mean-field family as a special case (where each $q_c$ is itself fully factored).

\paragraph{Computational cost.}
The dominant cost per sweep is the potential computation:
$O(|E| \cdot L \cdot |\alphabet|^4)$ for Match$\to$Match transitions
(contracting the $(|\alphabet|, |\alphabet|)$ pairwise marginal against the
$(|\alphabet|, |\alphabet|, |\alphabet|, |\alphabet|)$ WFST tensor).
The Felsenstein passes add $O(|E| \cdot L \cdot |\alphabet|^2)$, which is subdominant.


\subsection{Phylogenetic Hidden Markov Model (PhyloHMM)}

If the top-level indel rates in MixDom are low, and the ancestral presence/absence
fully specified by the MSA, the phylogenetic likelihood calculation and ancestral reconstruction
problems admit systematic approximation by a generalized Phylo-HMM,
yielding $O(L^2)$-complexity versions of the Forward and Forward-Backward algorithms.
This approach is described in Section~\ref{sec:partition-recon}.


